\clearpage
\SourceLeaf{479}
\pdfbookmark[1]{Posthumous Papers III. Special Perturbations of Pallas by Jupiter. First Calculation.}{pallas-III}
\fancyhead[L]{\small POSTHUMOUS PAPERS. PERTURBATIONS OF PALLAS}
\begin{center}
{\Large POSTHUMOUS PAPERS.\par}
\bigskip
{\large [III. SPECIAL PERTURBATIONS OF PALLAS BY JUPITER. FIRST CALCULATION.]\par}
\medskip
[OCTOBER—DECEMBER 1810.]
\end{center}
\par\Needspace{9\baselineskip}\bigskip
\pdfbookmark[2]{First Calculation, [1]}{pallas-III-fragment1}
\begin{center}[1.]\end{center}
\begingroup
\makeatletter\tagsleft@true\let\veqno\@@leqno\makeatother
\setlength{\lineskip}{4pt}\setlength{\lineskiplimit}{4pt}
\noindent\hspace*{\parindent}[The calculation concerns the perturbations of the elements and extends over the 3000-day period from 1803 June 30 to 1811 September 16. The period was divided into 60 intervals of 50 days each, so that the calculations had to be carried out for 61 instants in all. The calculation relies chiefly on formulas derived, with insignificant modifications, in the »Exposition d’une nouvelle méthode« etc. See Art. 20 of that work (page 464 of this volume).
\par
The following notation was used, taking the plane of Pallas’s orbit as the fundamental plane.
\begin{center}
\begin{tabular}{@{}r@{\ }p{.86\linewidth}@{}}
\(r={}\)&Radius vector of Pallas\\
\(v={}\)&heliocentric longitude of Pallas in its orbit\\
\(E={}\)&eccentric anomaly of Pallas\\
\(r'={}\)&Radius vector of Jupiter\\
\(w'={}\)&heliocentric longitude of Jupiter in its orbit\\
\(v'={}\)&longitude of Jupiter, projected onto the plane of Pallas’s orbit\\
\(\beta'={}\)&latitude of Jupiter above the plane of Pallas’s orbit\\
\(\rho={}\)&distance between Pallas and Jupiter\\
\(\mu={}\)&mass of Jupiter in units of the Sun’s mass\\
\(k={}\)&constant of the Theoria motus
\end{tabular}
\end{center}
\[
\begin{gathered}
\left.
\begin{aligned}
\begin{gathered}a,\ e,\ \varphi,\ \tilde\omega,\ p=a(1-ee),\\ n,\ \varepsilon,\ i,\ \AscendingNode\end{gathered}&=\text{elements of Pallas’s orbit}\\
a',\ e'\ \text{etc.}&=\text{elements of Jupiter’s orbit}
\end{aligned}
\right\}\quad\\
\text{with their usual meanings.}
\end{gathered}
\]
\par
For brevity we put
\[
\begin{aligned}
r'\cos\beta'\cos(v-v')&=\xi&\qquad\frac{1}{\rho^{3}}-\frac{1}{r'^{3}}&=\Delta\\
r'\cos\beta'\sin(v-v')&=\eta&&\\
r'\sin\beta'&=\zeta&\frac{\mu k\mathrm{d}t}{\BareRoot{p}}&=(1),
\end{aligned}
\tag*{1)}
\]
\SourceLeaf{480}which also gives \(rr+r'r'-2r\xi=\rho\rho\). The quantities \(T,V,W\), which here differ from the quantities denoted by the same symbols in the Exposition, Art. 15 (page 458), by the factors \(-\dfrac{a}{\cos\varphi}\), \(-\dfrac{ar}{\cos\varphi}\), \(+\dfrac{ar}{\cos\varphi}\), respectively, are calculated by the following formulas.
\[
T=\frac{(1)r}{\rho^{3}}-(1)\Delta\xi,\qquad V=(1)r\Delta\eta,\qquad W=(1)r\Delta\zeta,
\tag*{2)}
\]
These give the amounts of the instantaneous perturbations during the interval \(\mathrm{d}t\).
\[
\begin{aligned}
\mathrm{d}i&=\cos(v-\AscendingNode).W\\
\mathrm{d}\AscendingNode&=(2)\sin(v-\AscendingNode).W\\
\mathrm{d}n&=(4)\sin(v-\tilde\omega).T+\frac{(5)}{rr}V\\
\mathrm{d}\varphi&=(6)\sin(v-\tilde\omega).T+\frac{(6)(\cos(v-\tilde\omega)+\cos E)}{r}V\\
\mathrm{d}\tilde\omega&=A+(7)\cos(v-\tilde\omega).T\\
&\qquad+(8)\sin(\dot v-\tilde\omega).\left(1+\frac{p}{r}\right)V\\
\mathrm{d}\varepsilon&=A+(9)rT+(10)\cos(v-\tilde\omega).T\\
&\qquad+(11)\sin(v-\tilde\omega).\left(1+\frac{p}{r}\right)V,
\end{aligned}
\tag*{3)}
\]
Here, for brevity, the following notation is used.
\[
A=(3)\sin(v-\AscendingNode).W
\]
\[
\begin{aligned}
(2)&=\frac{1}{\sin i}&\qquad(7)&=\frac{p}{e}\\[3pt]
(3)&=\operatorname{tang}\tfrac{1}{2}i&(8)&=-\frac{1}{e}\\[3pt]
(4)&=3aen&(9)&=2\cos\varphi\\[3pt]
(5)&=3apn&(10)&=a\cos\varphi^{2}\operatorname{tang}\tfrac{1}{2}\varphi\\[3pt]
(6)&=-a\cos\varphi&(11)&=-\operatorname{tang}\tfrac{1}{2}\varphi.]
\end{aligned}
\]
\par\endgroup
